\documentclass{report}
\usepackage{amsmath,amssymb}
\setlength{\parindent}{0mm}
\setlength{\parskip}{1em}
\begin{document}
\begin{center}
\rule{6in}{1pt} \
{\large James P. Fairbanks \\
{\bf Eigenresidual Tolerances for Spectral Partitioning}}

855 W Peachtree St NW 1301 \\ Atlanta \\ GA 30308
\\
{\tt james.fairbanks@gatech.edu}\\
James P. Fairbanks\\
David A. Bader\end{center}

Spectral partitioning algorithms use eigenvectors and eigenvalues of the
graph Laplacian to partition graphs. This talks aims to provide insight
into the relationship between numerical accuracy and data mining quality
in the context of spectral partitioning. When the spectral gap of the
Laplacian is small, accurate approximations to the eigenvectors are hard
to compute. We provide a bound on the error between an approximate
eigenvector and the nearest element of an invariant subspace.
This shows quantitatively that accurate approximations to linear
combinations of eigenvectors are easier to compute when the spectral gap
is small. On a model problem where linear combinations of Laplacian
eigenvectors correctly partition the graph, the benefit of this shift in
perspective is large. Using this model problem, known as the Ring of
Cliques, and an analysis of sweep cut spectral partitioning, we derive
necessary and sufficient eigenresidual tolerances for iterative solvers
on this problem. This leads to a family of graphs where the necessary and
sufficient eigenresidual tolerance is $O(\frac{1}{\sqrt{n}})$. This
analysis provides guidance to practitioners about choosing an
eigenresidual tolerance for spectral partitioning applications on general
graphs. This informs decisions regarding error tolerances and reduces
iteration counts and computation time for practical applications such as
community detection and data clustering.


\end{document}
